\documentclass[11pt]{article}
\usepackage[margin=1in]{geometry}
\usepackage{amsmath,amssymb,amsthm}
\usepackage{xurl}
\usepackage{hyperref}
\sloppy
\let\oldus\_ \renewcommand{\_}{\oldus\allowbreak}
\newtheorem{theorem}{Theorem}
\newtheorem{lemma}[theorem]{Lemma}
\theoremstyle{definition}
\newtheorem{definition}[theorem]{Definition}
\newcommand{\R}{\mathbb{R}}

\title{Proof of conjecture 00000003287:\\ Gauss--Christoffel quadrature has its nodes at the zeros
of the orthogonal polynomial and maximal precision $2n-1$}
\author{Jackmeson1}
\date{2026-10-05}

\begin{document}
\maketitle

\section{The conjecture and its reading}

\textbf{Statement (English, verbatim).} Definition: Orthogonal polynomials: three-term recurrence
families of weight functions. Conjecture: Christoffel numerical integration: the nodes of Gauss
quadrature are polynomial zeros, with the algebraic precision of quadrature the maximum $2n-1$.
(Gauss quadrature 2n-1 maximal precision.) The Chinese text states the same.

This is the classical theorem on Gauss--Christoffel quadrature; we prove it (it is a known result,
not new). Quotations below are transliterated to ASCII. The Wikipedia article ``Gaussian
quadrature'' (retrieved in this session) states that an
$n$-point Gaussian rule ``is a quadrature rule constructed to yield an exact result for polynomials
of degree 2n - 1 or less by a suitable choice of the nodes xi and weights wi for i = 1, ..., n'',
and, in its proof of the fundamental theorem, ``Therefore pn has n distinct roots, all real, in the
interval from a to b.''

\textbf{Setting.} We work in the most general setting we could make faithful: a measure $\mu$ on
$\R$ such that
\begin{itemize}
\item (\texttt{FiniteMoments}) $x\mapsto x^k$ is $\mu$-integrable for every $k\in\mathbb{N}$
  (in particular $\mu$ is finite);
\item (\texttt{NotFinitelySupported}) $\mu(\R\setminus s)\neq 0$ for every finite set $s$,
  i.e.\ $\mu$ has infinitely many points of increase.
\end{itemize}
The second condition is the standard non-degeneracy condition: the Wikipedia article ``Orthogonal
polynomials'' says the inner product $\langle f,g\rangle=\int fg\,d\alpha$ ``is positive definite if
the function alpha has an infinite number of points of growth''. A weight function $w\ge 0$ on an
interval $[a,b]$ (or on $\R$) is the case $\mu=w(x)\,dx$ with $w=0$ off the interval; Lean lemma
\texttt{notFinitelySupported\_withDensity} proves that $w(x)\,dx$ satisfies the second condition
whenever $\int w\neq 0$. The moment condition is kept as a hypothesis.

\textbf{Definitions (all as in the Lean file).}
\begin{itemize}
\item $p$ is the \emph{monic orthogonal polynomial of degree $n$} (\texttt{IsMonicOrthogonal}) if
  $p$ is monic, $\deg p=n$, and $\int p\,q\,d\mu=0$ for every real polynomial $q$ with
  $\deg q<n$ (the zero polynomial has degree $-\infty$ and is included).
\item An \emph{$n$-node rule} is a pair of maps $x,w:\{0,\dots,n-1\}\to\R$ (nodes and weights;
  arbitrary, not necessarily distinct nodes, arbitrary real weights) and the functional
  $f\mapsto\sum_i w_i f(x_i)$. It is \emph{exact up to degree $d$} (\texttt{IsExactUpTo}) if
  $\int f\,d\mu=\sum_i w_i f(x_i)$ for every polynomial $f$ with $\deg f\le d$ (Lean's
  \texttt{natDegree}; the zero polynomial is included).
\item The \emph{algebraic precision} (\texttt{algebraicPrecision}) is
  $\sup\{d\in\mathbb{N} : \text{the rule is exact up to degree } d\}$ (Lean's \texttt{sSup} on
  $\mathbb{N}$; for the Gauss rule we prove this set has greatest element $2n-1$, so no
  convention for empty or unbounded sets is used).
\item The \emph{Christoffel numbers} of nodes $x$ (\texttt{christoffelWeight}) are
  $\lambda_i=\int \ell_i\,d\mu$, where $\ell_i=\prod_{j\neq i}(X-x_j)/(x_i-x_j)$ is the Lagrange
  basis polynomial (Mathlib's \texttt{Lagrange.basis}).
\end{itemize}

\textbf{Reading.} With $n\ge 1$ (natural-number subtraction $2n-1$ is then ordinary subtraction):
\begin{enumerate}
\item the orthogonal polynomial $p_n$ exists and is unique, and the ``three-term recurrence family''
  of the Definition is realized: $p_{n+2}=(X-a)p_{n+1}-b\,p_n$ with $b>0$;
\item $p_n$ has $n$ real zeros, all simple;
\item the Gauss rule (nodes the zeros $x_1<\dots<x_n$ of $p_n$, weights the Christoffel numbers)
  has positive weights, is exact up to degree $2n-1$ and not exact up to degree $2n$, so its
  algebraic precision is exactly $2n-1$;
\item ``the nodes of Gauss quadrature are polynomial zeros'': every $n$-node rule that is exact up
  to degree $2n-1$ satisfies $p_n=\prod_i (X-x_i)$, so its nodes are exactly the zeros of $p_n$;
\item ``the maximum $2n-1$'': every $n$-node rule (any nodes, any weights) exact up to degree $d$
  has $d\le 2n-1$.
\end{enumerate}

\section{Proof}

Write $I(f)=\int f\,d\mu$; by the moment condition every polynomial is integrable, so $I$ is a
linear functional on $\R[X]$ (Lean: \texttt{integrable\_eval}, \texttt{I}).

\begin{lemma}[\texttt{integral\_pos}]
If $f\neq 0$ and $f(x)\ge 0$ for all $x\in\R$, then $I(f)>0$.
\end{lemma}
\begin{proof}
$I(f)\ge 0$. If $I(f)=0$ then $f=0$ $\mu$-a.e., so $\mu$ is concentrated on the finite zero set of
$f$, contradicting \texttt{NotFinitelySupported}.
\end{proof}
In particular $I(r^2)=0$ forces $r=0$ (\texttt{eq\_zero\_of\_I\_mul\_self}).

\textbf{Existence and uniqueness} (\texttt{exists\_monicOrthogonal},
\texttt{monicOrthogonal\_unique}). Let $V$ be the polynomials of degree $<n$, identified with $\R^n$
by coefficients. The linear map $\Phi:V\to\R^n$, $\Phi(r)_i=I(rX^i)$, is injective: $\Phi(r)=0$
gives $I(rq)=0$ for all $q\in V$, hence $I(r^2)=0$ and $r=0$. So $\Phi$ is surjective, and choosing
$r$ with $\Phi(r)_i=-I(X^{n}X^i)$ gives the monic orthogonal $p_n=X^n+r$. If $p,p'$ are both monic
orthogonal of degree $n$, then $d=p-p'$ has degree $<n$ and $I(d^2)=I(pd)-I(p'd)=0$, so $p=p'$.

\textbf{Real simple zeros} (\texttt{roots\_card\_eq}, \texttt{roots\_nodup}). Key observation
(\texttt{not\_nonneg\_mul}): there is no $s\neq 0$ with $\deg s<n$ and $p_n s\ge 0$ on $\R$, since
then $0<I(p_ns)=0$. Write $p_n=P\,q$ with $P=\prod_{a}(X-a)$ over the real roots (with
multiplicity) and $q$ without real zeros (Mathlib \texttt{exists\_prod\_multiset\_X\_sub\_C\_mul}).
By the intermediate value theorem $q$ has constant sign (\texttt{const\_sign}). If $P$ had degree
$<n$, then $s=\pm P$ would give $p_ns=\pm q P^2\ge 0$, impossible; so $p_n$ has $n$ real roots with
multiplicity. If some root $a$ had multiplicity $\ge 2$, then $p_n=(X-a)^2s$ with $\deg s=n-2$ and
$p_ns=(X-a)^2s^2\ge 0$, impossible; so the roots are simple.

\textbf{Exactness up to $2n-1$} (\texttt{exact\_of\_roots}). Let $x_1,\dots,x_n$ be the zeros of
$p_n$ and $\deg f\le 2n-1$. Divide: $f=p_n\,d+r$ with $\deg r<n$ and $\deg d\le n-1$ (Mathlib
\texttt{modByMonic\_add\_div}). Then $I(p_nd)=0$, $r=\sum_i r(x_i)\ell_i$ (Lagrange interpolation,
\texttt{Lagrange.eq\_interpolate}) and $f(x_i)=r(x_i)$, so
$I(f)=I(r)=\sum_i r(x_i)\lambda_i=\sum_i\lambda_i f(x_i)$.

\textbf{Positive weights} (\texttt{weight\_pos}). For any rule with distinct nodes that is exact up
to degree $2n-1$, apply exactness to $\ell_i^2$ (degree $2n-2$): $w_i=\sum_j w_j\ell_i(x_j)^2=
I(\ell_i^2)>0$.

\textbf{Maximality} (\texttt{not\_exact\_two\_mul}, \texttt{precision\_le}). For any $n$-node rule
let $\omega=\prod_i(X-x_i)$, of degree $n$. Then $\omega^2$ has degree $2n$,
$I(\omega^2)>0$, and the rule gives $\sum_i w_i\omega(x_i)^2=0$. So no $n$-node rule is exact up to
degree $2n$, hence none is exact up to any $d\ge 2n$; i.e.\ every exactness degree is $\le 2n-1$.
For the Gauss rule the set of exactness degrees therefore has greatest element $2n-1$
(\texttt{precision\_eq}).

\textbf{Nodes are the zeros} (\texttt{nodes\_eq}). If an $n$-node rule is exact up to $2n-1$, then
for $\deg q<n$ the product $\omega q$ has degree $\le 2n-1$, so $I(\omega q)=\sum_i
w_i\omega(x_i)q(x_i)=0$. Thus $\omega$ is monic orthogonal of degree $n$, and by uniqueness
$\omega=p_n$.

\textbf{Three-term recurrence} (\texttt{three\_term\_recurrence}). Let $p_n,p_{n+1},p_{n+2}$ be
monic orthogonal. $Xp_{n+1}-p_{n+2}$ has degree $\le n+1$; subtracting $a\,p_{n+1}$ and then
$b\,p_n$ (with $a,b$ the appropriate top coefficients, \texttt{degree\_sub\_lead}) leaves
$v$ of degree $<n$ with $I(vX^i)=0$ for $i<n$ (using $I(p_{n+1}X^{i+1})=0$), so $v=0$:
$p_{n+2}=(X-a)p_{n+1}-b\,p_n$. Finally $b\,I(p_n^2)=I(p_{n+1}\cdot Xp_n)=I(p_{n+1}^2)>0$, so
$b>0$.

\section{Lean formalization and axiom audit}

File \texttt{lean/Conjecture3287/Basic.lean} (namespace \texttt{C3287}; only import
\texttt{Mathlib}). The main theorem is
\begin{verbatim}
theorem gauss_christoffel (hmom : FiniteMoments mu)
    (hsupp : NotFinitelySupported mu) {n : Nat} (hn : 0 < n) :
    (exists! p : R[X], IsMonicOrthogonal mu n p) /\
    forall p : R[X], IsMonicOrthogonal mu n p ->
      (Multiset.card p.roots = n /\ p.roots.Nodup) /\
      (exists x : Fin n -> R, StrictMono x /\
        (forall a : R, p.IsRoot a <-> a in Set.range x) /\
        (forall i, 0 < christoffelWeight mu x i) /\
        IsExactUpTo mu x (christoffelWeight mu x) (2 * n - 1) /\
        not IsExactUpTo mu x (christoffelWeight mu x) (2 * n) /\
        algebraicPrecision mu x (christoffelWeight mu x) = 2 * n - 1) /\
      (forall x w : Fin n -> R, IsExactUpTo mu x w (2 * n - 1) ->
        p = prod i, (X - C (x i))) /\
      (forall x w : Fin n -> R, forall d, IsExactUpTo mu x w d -> d <= 2 * n - 1)
\end{verbatim}
(ASCII transliteration of the Lean statement). Further top-level results:
\texttt{three\_term\_recurrence} (item 1) and \texttt{notFinitelySupported\_withDensity} (weight
functions $w(x)\,dx$ are admissible).

\textbf{Axiom audit.} \texttt{lake build} succeeds; \texttt{\#print axioms} for
\texttt{C3287.gauss\_christoffel}, \texttt{C3287.three\_term\_recurrence} and
\texttt{C3287.notFinitelySupported\_withDensity} reports only \texttt{propext},
\texttt{Classical.choice}, \texttt{Quot.sound}. No \texttt{sorry}, no \texttt{native\_decide}, no
added axioms.

\section{Sources}
\begin{itemize}
\item Wikipedia, ``Gaussian quadrature'', \url{https://en.wikipedia.org/wiki/Gaussian_quadrature}
  (quoted above).
\item Wikipedia, ``Orthogonal polynomials'', \url{https://en.wikipedia.org/wiki/Orthogonal_polynomials}
  (quoted above).
\end{itemize}

\end{document}
